\documentclass[10pt,letterpaper]{article}
\usepackage[utf8]{inputenc}
\usepackage[spanish,es-tabla]{babel}
\usepackage{geometry}
\geometry{margin=0.68in, top=0.85in, bottom=0.78in, headheight=36pt, headsep=14pt}
\usepackage{amsmath,amssymb}
\usepackage{xcolor}
\usepackage{graphicx}
\usepackage{tabularx}
\usepackage{booktabs}
\usepackage{colortbl}
\usepackage{array}
\usepackage{enumitem}
\usepackage{fancyhdr}
\usepackage{titlesec}
\usepackage{lastpage}
\usepackage{setspace}
\usepackage{listings}

% Tipografía institucional Icesi (Sans-Serif Helvetica)
\usepackage{helvet}
\renewcommand{\familydefault}{\sfdefault}

% Interlineado aumentado un 30% (setstretch 1.38) para mayor holgura y lectura óptima
\setstretch{1.38}

% Altura de celdas en tablas
\renewcommand{\arraystretch}{1.25}

% Paleta de Colores Oficial de los Programas Icesi
\definecolor{icesiBlue}{HTML}{5454E9}       % Azul vibrante institucional
\definecolor{icesiOrange}{HTML}{F0681A}     % Naranja acento
\definecolor{icesiGreen}{HTML}{4CB978}      % Verde acento
\definecolor{icesiPurple}{HTML}{865CF0}     % Púrpura acento
\definecolor{icesiDark}{HTML}{1A1A1A}       % Texto principal oscuro
\definecolor{icesiMuted}{HTML}{6B7280}      % Gris metadatos
\definecolor{icesiBorder}{HTML}{D1D5DB}     % Líneas y bordes
\definecolor{icesiRowLine}{HTML}{E5E7EB}    % Líneas de tabla
\definecolor{icesiLight}{HTML}{F9FAFB}      % Fondo suave
\definecolor{codeBackground}{HTML}{F3F4F6}  % Fondo código

% Configuración de Listings
\lstdefinestyle{bashstyle}{
    backgroundcolor=\color{codeBackground},
    commentstyle=\color{icesiMuted}\itshape,
    keywordstyle=\color{icesiBlue}\bfseries,
    basicstyle=\ttfamily\scriptsize,
    breaklines=true,
    frame=single,
    rulecolor=\color{icesiBorder},
    language=bash,
    aboveskip=2pt,
    belowskip=2pt
}

\lstdefinestyle{codeStyle}{
    backgroundcolor=\color{codeBackground},
    commentstyle=\color{icesiMuted}\itshape,
    keywordstyle=\color{icesiBlue}\bfseries,
    stringstyle=\color{icesiGreen},
    basicstyle=\ttfamily\scriptsize,
    breaklines=true,
    frame=single,
    rulecolor=\color{icesiBorder},
    aboveskip=2pt,
    belowskip=2pt
}

% =========================================================================
% ENCABEZADO Y PIE DE PÁGINA ESTILO PROGRAMA DE CURSO ICESI
% =========================================================================
\pagestyle{fancy}
\fancyhf{}
\fancyhead[L]{%
  \begin{minipage}[c]{2.8cm}
    \includegraphics[height=0.58cm]{logo-icesi.png}
  \end{minipage}%
  \hspace{0.15cm}%
  {\color{icesiBorder}\vrule width 0.8pt height 0.60cm}%
  \hspace{0.2cm}%
  \begin{minipage}[c]{10.5cm}
    \fontsize{7.5}{9.5}\selectfont\textbf{\color{icesiDark}UNIVERSIDAD ICESI} \\
    \fontsize{7}{9}\selectfont\textbf{\color{icesiDark}PROGRAMA DE CURSO / TALLER EVALUATIVO: INGENIERÍA DE SOFTWARE V} \\
    \fontsize{6.5}{8.5}\selectfont\color{icesiMuted}Período: 202620 $\bullet$ NRC: 11005 $\bullet$ Duración Máxima: 2 Horas
  \end{minipage}%
}
\fancyhead[R]{\fontsize{7.5}{9.5}\selectfont\color{icesiMuted}Pág. \thepage\ de \pageref{LastPage}}

% Barra horizontal multicolor característica del programa de curso
\renewcommand{\headrule}{%
  \vspace{0.08cm}%
  \noindent%
  \textcolor{icesiBlue}{\rule{0.60\textwidth}{1.8pt}}%
  \textcolor{icesiOrange}{\rule{0.18\textwidth}{1.8pt}}%
  \textcolor{icesiGreen}{\rule{0.11\textwidth}{1.8pt}}%
  \textcolor{icesiPurple}{\rule{0.11\textwidth}{1.8pt}}%
}
\renewcommand{\footrulewidth}{0.4pt}
\fancyfoot[L]{\fontsize{7.5}{9.5}\selectfont\color{icesiMuted}Ingeniería de Software V (09776 - TIC) $\bullet$ Taller Pruebas E2E con Playwright \& CI/CD}
\fancyfoot[R]{\fontsize{7.5}{9.5}\selectfont\color{icesiMuted}Universidad Icesi $\bullet$ DCSI}

% =========================================================================
% COMANDOS DE SECCIONES CON BADGE ESTILO PROGRAMA ICESI
% =========================================================================
\newcommand{\sectionBadge}[1]{%
  \refstepcounter{section}%
  \vspace{0.16cm}%
  \noindent%
  \colorbox{icesiBlue}{\color{white}\rule[-0.08cm]{0pt}{0.42cm}\hspace{0.16cm}\textbf{\normalsize \thesection}\hspace{0.16cm}}%
  \hspace{0.18cm}%
  {\fontsize{9.5}{12}\selectfont\textbf{\color{icesiDark}#1}}\\[-0.10cm]%
  \noindent{\color{icesiDark}\rule{\textwidth}{0.8pt}}\vspace{0.06cm}%
}

\titleformat{\subsection}
  {\color{icesiBlue}\normalfont\normalsize\bfseries}
  {\thesubsection}{0.4em}{}
\titlespacing*{\subsection}{0pt}{0.4ex plus 0.2ex minus .1ex}{0.2ex plus .1ex}

\newcolumntype{C}[1]{>{\centering\arraybackslash}p{#1}}
\newcolumntype{L}[1]{>{\raggedright\arraybackslash}p{#1}}
\newcolumntype{Y}{>{\centering\arraybackslash}X}

\begin{document}

% =========================================================================
% TÍTULO Y PRESENTACIÓN INICIAL (PÁGINA 1)
% =========================================================================
\noindent
{\fontsize{8.5}{11}\selectfont\textbf{\color{icesiBlue}TALLER EVALUATIVO: PRUEBAS END-TO-END (E2E) CON PLAYWRIGHT \& CI/CD (TIEMPO: 2 HORAS)}}\\[0.06cm]
{\fontsize{15}{18}\selectfont\textbf{Ingeniería de Software V}}\\[0.10cm]
{\color{icesiRowLine}\hrule height 0.6pt}\vspace{0.10cm}

\noindent
{\color{icesiBlue}\textbf{\rule{2.5pt}{8.5pt}}}~\textbf{\fontsize{8.5}{11}\selectfont FACULTAD DE INGENIERÍA, DISEÑO Y CIENCIAS APLICADAS} \hfill
{\color{icesiOrange}\textbf{\rule{2.5pt}{8.5pt}}}~\textbf{\fontsize{8.5}{11}\selectfont DEPARTAMENTO DE COMPUTACIÓN Y SISTEMAS INTELIGENTES}\\[0.06cm]
{\color{icesiRowLine}\hrule height 0.6pt}\vspace{0.10cm}

% =========================================================================
% TABLA DE INFORMACIÓN Y METADATOS (ESTILO PROGRAMA DE CURSO)
% =========================================================================
\noindent
{\color{icesiBlue}\textbf{\rule{2.5pt}{8.5pt}}}~\textbf{\fontsize{8.5}{11}\selectfont INFORMACIÓN DE LA ACTIVIDAD} \hfill {\fontsize{7.5}{9}\selectfont\color{icesiMuted}\textbf{IDENTIFICACIÓN}}\\[-0.22cm]
{\color{icesiBorder}\hrule height 0.5pt}\vspace{0.05cm}

\noindent
\small
\begin{tabularx}{\textwidth}{>{\bfseries\color{icesiDark}\fontsize{7.5}{9.5}\selectfont}L{2.5cm}X!{\color{icesiBlue}\vrule width 0.8pt}>{\bfseries\color{icesiDark}\fontsize{7.5}{9.5}\selectfont}L{2.5cm}X}
CÓDIGO: & 09776 - TIC & PERÍODO: & 202620 (Semestre 2026-2) \\ \hline
\rowcolor{icesiLight}
NRC / GRUPO: & 11005 / Grupo 001 & PONDERACIÓN: & 15.0\% (Taller Evaluativo) $\bullet$ Duración: 2 Horas \\ \hline
POLÍTICA IAG: & Nivel 3: Colaboración Asistida & CALIFICACIÓN: & \rule{1.5cm}{0.4pt} / 5.0 \\ \hline
\rowcolor{icesiLight}
ESTUDIANTE 1: & \rule{4.5cm}{0.4pt} & CÓDIGO: & \rule{2.2cm}{0.4pt} \\ \hline
ESTUDIANTE 2: & \rule{4.5cm}{0.4pt} & CÓDIGO: & \rule{2.2cm}{0.4pt} \\ \hline
\end{tabularx}

\vspace{0.10cm}

% =========================================================================
% SECCIÓN 1: CONTEXTO Y DESCRIPCIÓN DEL RETO
% =========================================================================
\sectionBadge{Contexto del Problema \& Aseguramiento de Calidad E2E}

\subsection{Descripción del Desafío de Ingeniería de Software}
La organización \textbf{Icesi Cloud Solutions} opera una plataforma de comercio desacoplada compuesta por un Backend REST en \textbf{Spring Boot 3 (Java 17)} y un cliente web en \textbf{React 18 + Vite + TypeScript}.

Aunque el equipo cuenta con pruebas unitarias en Maven que reportan 100\% de éxito, despliegues recientes a producción generaron fallas críticas de usuario: regresiones en formularios asíncronos, inconsistencias entre navegadores (Chromium, Firefox y WebKit) y congelamiento de la interfaz ante caídas de red.

El objetivo de este taller es estructurar una suite profesional de pruebas automatizadas \textbf{End-to-End (E2E)} con \textbf{Playwright}, implementando el patrón \textbf{Page Object Model (POM)}, localizadores resilientes basados en accesibilidad (\textbf{W3C / ARIA}), simulación de contingencias (\textbf{Network Mocking}) e integración continua en \textbf{GitHub Actions} con publicación automatizada de reportes y trazas.

\vspace{0.08cm}
\noindent
\small
\begin{tabularx}{\textwidth}{|L{2.8cm}|X|L{3.5cm}|}
\hline
\rowcolor{icesiLight}
\textbf{Componente} & \textbf{Tecnología \& Stack} & \textbf{Puerto / Endpoint} \\ \hline
\textbf{Backend REST} & Spring Boot 3.2, Spring Security (JWT), H2 & Puerto \texttt{8080} (\texttt{/api/v1/orders}) \\ \hline
\textbf{Frontend SPA} & React 18, TypeScript, Vite, TailwindCSS & Puerto \texttt{3000} (\texttt{http://localhost:3000}) \\ \hline
\textbf{Suite E2E} & Playwright Test Runner v1.42+, TypeScript & Motores: Chromium, Firefox, WebKit \\ \hline
\textbf{CI/CD Automation} & GitHub Actions (\texttt{e2e-playwright.yml}) & Headless Execution + HTML Report \\ \hline
\end{tabularx}

\newpage

% =========================================================================
% PÁGINA 2: PROCEDIMIENTO EXPERIMENTAL (FASES 1 Y 2)
% =========================================================================
\sectionBadge{Procedimiento Experimental \& Tareas a Desarrollar}

\subsection{Fase 1: Configuración de Playwright, Tooling y Multi-Navegador (25 min)}
En el directorio raíz del proyecto de pruebas:
\begin{enumerate}[leftmargin=*,label=\textbf{\alph*)},itemsep=1.5pt,topsep=1.5pt]
  \item Inicialice la estructura de pruebas asegurando soporte TypeScript y navegadores oficiales:
  \begin{lstlisting}[style=bashstyle]
npm init playwright@latest --yes -- --typescript --browsers=all --github-actions
  \end{lstlisting}
  \item Configure el archivo central \texttt{playwright.config.ts} cumpliendo con los siguientes parámetros:
  \begin{itemize}[itemsep=1pt]
    \item \texttt{baseURL: 'http://localhost:3000'} con timeout por prueba de 30 segundos.
    \item Política de reintentos en CI: \texttt{retries: process.env.CI ? 2 : 0} para mitigar latencias de red.
    \item Diagnóstico: \texttt{trace: 'on-first-retry'}, capturas automáticas (\texttt{screenshot: 'only-on-failure'}) y grabación (\texttt{video: 'retain-on-failure'}).
    \item Proyectos de ejecución multi-motor: \texttt{chromium}, \texttt{firefox} y \texttt{webkit}.
  \end{itemize}
  \item Incorpore en \texttt{package.json} los scripts: \texttt{"test:e2e": "playwright test"}, \texttt{"test:e2e:ui": "playwright test --ui"}, y \texttt{"test:e2e:report": "playwright show-report"}.
\end{enumerate}

\subsection{Fase 2: Patrón Page Object Model (POM) y Flujos Críticos de Negocio (45 min)}
Desacople la interacción de la interfaz de la lógica de aserción creando clases modulares en \texttt{e2e/pages/}:
\begin{enumerate}[leftmargin=*,label=\textbf{\alph*)},itemsep=1.5pt,topsep=1.5pt]
  \item \textbf{Clase \texttt{LoginPage.ts}:}
    \begin{itemize}[itemsep=1pt]
      \item Encapsule la navegación a la vista (\texttt{goto('/login')}).
      \item Defina localizadores semánticos basados en ARIA: \texttt{page.getByRole('textbox', \{ name: /usuario|correo/i \})}, \texttt{page.getByLabel(/contraseña/i)}, y \texttt{page.getByRole('button', \{ name: /iniciar sesión|ingresar/i \})}.
      \item Implemente el método \texttt{async login(user, pass)} y el localizador de alertas: \texttt{page.getByRole('alert')}.
    \end{itemize}
  \item \textbf{Clase \texttt{CatalogPage.ts}:}
    \begin{itemize}[itemsep=1pt]
      \item Localizador para el buscador: \texttt{page.getByPlaceholder(/buscar productos/i)}.
      \item Método para filtrar y seleccionar un producto: \texttt{page.getByRole('article').filter(\{ hasText: name \})}.
      \item Método para agregar al carrito y consultar la cantidad del badge: \texttt{getCartBadgeCount()}.
    \end{itemize}
  \item \textbf{Clase \texttt{CheckoutPage.ts}:}
    \begin{itemize}[itemsep=1pt]
      \item Validación de subtotales y diligenciamiento de entrega.
      \item Confirmación de compra con \texttt{page.getByRole('button', \{ name: /confirmar orden/i \})}.
      \item Verificación del modal de confirmación con el identificador de orden generado.
    \end{itemize}
  \item \textbf{Suite de Pruebas (\texttt{tests/e2e/order-lifecycle.spec.ts}):}
    \begin{itemize}[itemsep=1pt]
      \item \textit{Test 1 (Camino Feliz):} Login $\rightarrow$ Catálogo $\rightarrow$ Agregar al carrito $\rightarrow$ Checkout $\rightarrow$ Aserción web-first de compra exitosa (\texttt{expect(modal).toBeVisible()}).
      \item \textit{Test 2 (Credenciales Inválidas):} Intento de acceso erróneo $\rightarrow$ Aserción de alerta visible sin redirección de URL (\texttt{expect(page).toHaveURL(/login/)}).
      \item \textit{Test 3 (Restricciones):} Comprobar que el botón de pago esté deshabilitado (\texttt{toBeDisabled()}) con carrito vacío.
    \end{itemize}
\end{enumerate}

\newpage

% =========================================================================
% PÁGINA 3: FASE 3, FASE 4 Y CUESTIONARIO TÉCNICO
% =========================================================================
\subsection{Fase 3: Interceptación de Red (Network Mocking) y Resiliencia (25 min)}
En el archivo \texttt{tests/e2e/network-resilience.spec.ts}, valide la tolerancia a fallos del frontend:
\begin{enumerate}[leftmargin=*,label=\textbf{\alph*)},itemsep=1.5pt,topsep=1.5pt]
  \item \textbf{Simulación de Error 500 en Checkout:} Interceptar con \texttt{page.route()} la petición \texttt{POST **/api/v1/orders} forzando código HTTP 500:
    \begin{lstlisting}[style=codeStyle]
await page.route('**/api/v1/orders', async (route) => {
  await route.fulfill({ status: 500, contentType: 'application/json',
    body: JSON.stringify({ message: 'Error interno en pasarela de pagos' }) });
});
    \end{lstlisting}
    Valide mediante aserciones web-first que la interfaz muestre el diálogo de reintento sin congelarse.
  \item \textbf{Simulación de Latencia de Red:} Intercepte \texttt{GET **/api/v1/products} agregando un retraso artificial de 3000 ms. Verifique que durante la espera se despliegue el indicador visual (\textit{spinner} o \textit{skeleton}) y se oculte tras completar la carga.
  \item \textbf{Diagnóstico con Trace Viewer:} Ejecute \texttt{npx playwright test --trace on} e inspeccione con \texttt{npx playwright show-trace} eventos DOM, llamadas HTTP y capturas paso a paso.
\end{enumerate}

\subsection{Fase 4: Pipeline Declarativo en GitHub Actions (15 min)}
Configure el flujo de integración continua en \texttt{.github/workflows/e2e-playwright.yml}:
\begin{enumerate}[leftmargin=*,label=\textbf{\alph*)},itemsep=1.5pt,topsep=1.5pt]
  \item Disparador ante eventos \texttt{push} y \texttt{pull\_request} hacia ramas de entrega.
  \item Job sobre \texttt{ubuntu-latest} con instalación de Node 20 y navegadores con dependencias:
  \begin{lstlisting}[style=bashstyle]
- name: Install dependencies & Playwright Browsers
  run: npm ci && npx playwright install --with-deps
- name: Run Playwright tests
  run: npx playwright test
  \end{lstlisting}
  \item Publicación del reporte HTML como artefacto descargable con retención de 30 días (\texttt{if: always()}):
  \begin{lstlisting}[style=codeStyle]
- uses: actions/upload-artifact@v4
  if: always()
  with: { name: playwright-report, path: playwright-report/ }
  \end{lstlisting}
\end{enumerate}

\vspace{0.06cm}

% =========================================================================
% SECCIÓN 3: PREGUNTAS DE ANÁLISIS TÉCNICO
% =========================================================================
\sectionBadge{Cuestionario de Análisis Conceptual \& Bitácora IAG (10 min)}

Responda brevemente a las siguientes preguntas con base en la práctica de laboratorio:
\begin{enumerate}[leftmargin=*,label=\textbf{\arabic*.},itemsep=2pt,topsep=1.5pt]
  \item \textbf{Localizadores Semánticos vs Frágiles:} ¿Por qué se desaconseja XPath o CSS acoplados al DOM en favor de métodos ARIA (\texttt{getByRole})? ¿Qué ventajas otorga ante refactorizaciones y accesibilidad WCAG?
  \item \textbf{Aislamiento de Contextos (\textit{BrowserContext}):} Compare la apertura de contextos en una sola instancia de Chromium frente a levantar procesos independientes (Selenium) en velocidad, memoria y aislamiento.
  \item \textbf{Network Mocking vs E2E Puros:} ¿Cuándo conviene interceptar llamadas con \texttt{page.route()} y cuándo es indispensable probar contra microservicios y bases de datos reales? ¿Qué riesgos acarrea abusar de mocks?
  \item \textbf{Pruebas Intermitentes (\textit{Flaky Tests}):} Explique cómo el \textit{auto-waiting} de Playwright elimina pausas fijas (\texttt{sleep}). ¿Es la política de \texttt{retries} en CI una solución definitiva o paliativa?
  \item \textbf{Bitácora y Uso Ético de IAG (Nivel 3):} ¿Qué herramientas de IA utilizó para estructurar clases POM o fixtures? ¿Qué alucinaciones o errores detectó y cómo los corrigió técnicamente?
\end{enumerate}

\newpage

% =========================================================================
% PÁGINA 4: RÚBRICA MODULARIZADA DE EVALUACIÓN Y ENTREGA
% =========================================================================
\sectionBadge{Rúbrica de Evaluación Modularizada por Punto (Escala 0.0 a 5.0)}

\vspace{0.08cm}
\noindent
\setlength{\tabcolsep}{3.5pt}
\fontsize{7.5}{9.5}\selectfont
\renewcommand{\arraystretch}{1.18}
\begin{tabularx}{\textwidth}{|L{2.5cm}|C{0.6cm}|X|X|X|}
\hline
\rowcolor{icesiLight}
\textbf{Punto Entregable} & \textbf{Pts} & \textbf{Excelente (4.5 -- 5.0)} & \textbf{Aceptable (3.0 -- 4.4)} & \textbf{Insuficiente (0.0 -- 2.9)} \\ \hline
\textbf{2.1. Configuración Playwright \& Tooling} & \textbf{0.8} & Configuración completa en \texttt{playwright.config.ts} multi-navegador (Chromium, Firefox, WebKit), trazas en fallo y scripts npm limpios. & Configuración básica en un solo motor de renderizado o sin políticas de diagnóstico en fallos. & Proyecto no inicializa correctamente o faltan dependencias del sistema operativo. \\ \hline
\textbf{2.2. Page Object Model \& Flujos Críticos} & \textbf{1.5} & Clases POM limpias en \texttt{pages/}, uso exclusivo de localizadores semánticos (\texttt{getByRole}), flujo completo de compra y aserciones web-first. & Flujos funcionales pero mezcla selectores directos en los tests o usa selectores CSS frágiles. & Pruebas incompletas, ausencia de POM o uso abusivo de pausas fijas (\texttt{waitForTimeout}). \\ \hline
\textbf{2.3. Network Mocking \& Resiliencia} & \textbf{1.0} & Interceptación elegante con \texttt{page.route()} simulando error 500 y latencia, verificación de retroalimentación en UI y uso de Trace Viewer. & Simulación básica de mock sin validación de estados de carga ni manejo de errores en UI. & No implementa interceptación de red o pruebas fallan por mal manejo de promesas. \\ \hline
\textbf{2.4. CI/CD Pipeline en GitHub Actions} & \textbf{1.0} & Workflow automatizado, instalación de navegadores con dependencias (\texttt{--with-deps}), ejecución headless y reporte HTML con \texttt{if: always()}. & Pipeline funcional pero sin persistencia adecuada de artefactos o falla intermitentemente. & Pipeline inexistente, sintaxis YAML inválida o permanentemente en estado fallido (\textit{Failing}). \\ \hline
\textbf{3.1. Cuestionario \& Bitácora IAG} & \textbf{0.7} & Respuestas técnicas profundas a las 5 preguntas del cuestionario con análisis riguroso y bitácora ética y crítica de uso de IAG. & Respuestas breves con justificación superficial o bitácora de prompts incompleta. & Respuestas erróneas conceptualmente o ausencia de la sección de análisis. \\ \hline
\rowcolor{icesiLight}
\textbf{CALIFICACIÓN TOTAL} & \textbf{5.0} & \multicolumn{3}{c|}{\textbf{Escala Institucional Universidad Icesi: Aprobatoria $\ge$ 3.0 $\bullet$ Duración Máxima: 2 Horas}} \\ \hline
\end{tabularx}

\vspace{0.20cm}

% =========================================================================
% SECCIÓN DE ENTREGABLES
% =========================================================================
\noindent
{\color{icesiBlue}\textbf{\rule{2.5pt}{8.5pt}}}~\textbf{\fontsize{8.5}{11}\selectfont INSTRUCCIONES DE ENTREGA} \hfill {\fontsize{7.5}{9}\selectfont\color{icesiMuted}\textbf{MODALIDAD}}\\[-0.22cm]
{\color{icesiBorder}\hrule height 0.5pt}\vspace{0.08cm}

\begin{enumerate}[leftmargin=*,itemsep=2pt,topsep=2pt]
  \item \textbf{Repositorio en GitHub:} Subir el código fuente completo con \texttt{pages/}, \texttt{tests/}, configuración de Playwright y workflow en \texttt{.github/workflows/e2e-playwright.yml}. La rama evaluada debe exhibir el último commit con el pipeline en \textbf{Verde (\textit{Passing})}.
  \item \textbf{Documento de Resultados (PDF):} Adjuntar en la plataforma institucional este documento debidamente diligenciado con:
    \begin{itemize}[itemsep=1pt]
      \item Nombres y códigos de los dos integrantes del equipo.
      \item Enlace web al repositorio en GitHub.
      \item Captura de pantalla del reporte HTML de Playwright (\texttt{playwright-report/index.html}) con el 100\% de las pruebas pasando en los tres motores.
      \item Captura de pantalla de la ejecución exitosa de la acción en GitHub Actions con el artefacto descargado.
      \item Respuestas detalladas al cuestionario conceptual y bitácora ética de uso de IAG de la Sección 3.
    \end{itemize}
\end{enumerate}

\end{document}
